\section{Розділ~\ref{sec:dofa}. Навчання з \nt{DOPA}}

Механізми синапсу (порції, детектор збігів, кальцій, вікна пластичності) і поведінку дофаміну як помилки передбачення виміряно прямо; те, що правила ведуть до градієнта і чого коштує широкомовлення, --- теореми; підсумок навчання вибору --- розрахунок за моделлю книги; схема, що обчислює помилку, --- гіпотеза.

{\footnotesize
\setlength{\tabcolsep}{3pt}\renewcommand{\arraystretch}{1.15}
\begin{longtable}{>{\raggedright}p{0.5cm}>{\raggedright}p{4.0cm}>{\raggedright}p{5.6cm}>{\raggedright}p{2.3cm}>{\raggedright\arraybackslash}p{1.7cm}}
\toprule
№ & Твердження & Дослід і що виміряно & Джерело & Статус \\
\midrule\endfirsthead
\toprule
№ & Твердження & Дослід і що виміряно & Джерело & Статус \\
\midrule\endhead
1 & Зворотного поширення помилки в тій формі, що в штучних мережах, у мозку не знайдено & Прямого досліду немає: це твердження про відсутність. Огляди розбирають, чого вимагає алгоритм (зворотні зв'язки, що повторюють прямі, окрема фаза) і які біологічні наближення до нього запропоновано. & \cite{Lillicrap2020, WhittingtonBogacz2019} & гіпотеза \\
2 & Вага розкладається на кількість місць викиду, імовірність викиду і відповідь на порцію: $w=N_s\,p\,A_1$~\eqref{eq:quantal} & Нервово-м'язовий синапс жаби за зниженого викиду: відповідь отримувача коливається сходинками, кратними спонтанній мініатюрній відповіді на одну порцію; кількість порцій на імпульс підлягає закону Пуассона. & \cite{delCastillo1954} & виміряно \\
3 & Рецептор отримувача --- детектор збігів; кальцій запускає зміну ваги & Патч-клемп нейронів миші: іони Mg$^{2+}$ закупорюють канал, відкритий глутаматом, за потенціалу спокою і виходять за деполяризації. Зрізи гіпокампа: хелатор кальцію в отримувачі блокує довготривале посилення, а вивільнення кальцію світлом усередині клітини саме посилює передачу. & \cite{Nowak1984, Malenka1988calcium} & виміряно \\
4 & Рішення виконує будь-який бік: росте $A_1$ в отримувача, змінюється $p$ у відправника & Глутамат, вивільнений світлом над одним шипиком, спричиняє швидкий і стійкий ріст шипика і струму через його рецептори; посилення супроводжується вбудовуванням AMPA-рецепторів. Деполяризація отримувача на певний час пригнічує викид у відправника; ефект переносять ендоканабіноїди, що йдуть назад через щілину (показано на \nt{GABA}-входах). Короткочасне виснаження залежить від імовірності викиду у відправника. & \cite{Matsuzaki2004, Malinow2002, WilsonNicoll2001, Tsodyks1997} & виміряно \\
5 & Синапс посилюється, коли відправник і отримувач активні разом~\eqref{eq:hebb} & Кролик під наркозом: після коротких серій імпульсів вхідним шляхом гіпокампа відповідь посилена на час від $30$~хв до $10$~год. У зрізі гіпокампа імпульси отримувача посилюють лише ті синапси, які в ту саму мить активні. & \cite{Bliss1973LTP, Kelso1986hebb} & виміряно \\
6 & Правило~\eqref{eq:oja} знаходить головний власний вектор кореляцій входів (твердження~\ref{prop:oja}) & Теорема; доведення в розділі. Досліду, де нейрон вивчив головну компоненту своїх входів, немає; механізм обмеження росту ваг у синапсі не встановлено. & \cite{Oja1982} & теорія \\
7 & Гебб за часом: вікно близько $\pm20\ms$ (рис.~\ref{fig:stdp}); з повторюваних послідовностей виростають ланцюжки & Пари нейронів гіпокампа в культурі: імпульс отримувача в межах $20\ms$ після імпульсу відправника дає стійке посилення, у межах $20\ms$ до --- послаблення; обидва залежать від NMDA-рецепторів. Відношення $A_-=0.6A_+$ на рисунку --- ілюстрація. Що так у мережі виростають ланцюжки, прямо не виміряно. & \cite{BiPoo1998} & виміряно \\
8 & Слід~\eqref{eq:threefactor}: дофамін, що надійшов через $1$--$2\,\text{с}$ після спільної активності, посилює зв'язок, пізніше --- ні (твердження~\ref{prop:threefactor}) & Зрізи стріатума, роздільна оптогенетична стимуляція глутаматних і дофамінових входів: шипик росте, лише якщо дофамін надходить через $0.3$--$2\,\text{с}$ після глутаматного входу; вікно задається швидким підйомом і розпадом цАМФ. У корі спарювання лишає окремі сліди для посилення і послаблення, які моноамінові рецептори перетворюють на довготривалі зміни у вікні порядку секунд. & \cite{Yagishita2014, He2015eligibility} & виміряно \\
9 & Вікна завдовжки в секунди бувають і без дофаміну: третій сигнал --- сильне збудження отримувача & Жива миша, поле CA1: одна плато-подія в отримувачі посилює входи, що надійшли за секунди до і після неї, і за один прохід створює поле місця. & \cite{Bittner2017} & виміряно \\
10 & Середній крок правила трьох множників іде за градієнтом винагороди~\eqref{eq:perturbation} (твердження~\ref{prop:gradient}) & Теорема про градієнт для навчання за випадковими відхиленнями; у розділі виведена для малого шуму. & \cite{Williams1992} & теорія \\
11 & Шум --- це пошук: мережа навчається на своїх випадкових відхиленнях & Молоді амадини: вимкнення ядра, що виходить із кола базальних гангліїв, різко і зворотно прибирає мінливість пісні. Дорослі амадини: завада подається на виконання складу з висотою по один бік від середньої, і птахи швидко зсувають висоту в інший бік --- навчаються за власним розкидом. & \cite{Olveczky2005, TumerBrainard2007} & виміряно \\
12 & Вибір з двох навчається при $\tau_e=1\,\text{с}$ і $\delta=R-\bar R$; не навчається при $\tau_e=50\ms$; без віднімання ваги ростуть без меж, а за великої швидкості навчання мережа може закріпити гірший вибір & \texttt{models/dofa\_learning.py}, $\eta=0.05$, $500$ спроб, $6$ зерен. $\tau_e=1\,\text{с}$: частка виборів $A$ $0.65$--$0.79\to0.96$--$1.00$, ваги $0.85$--$1.33$. $\tau_e=50\ms$: $0.38$--$0.55$, ваги не змінюються. $\delta=R$: вага переможця $860$--$1190$. $\delta=R$, $\eta=0.5$, $3$ зерна: одне закріпилося на $B$ ($0.04\to0$). Скрипт друкує всі чотири випадки; перерахунок збігся з розділом. & \texttt{models/\allowbreak dofa\_\allowbreak learning.py} & модель \\
13 & Час навчання за одним спільним сигналом росте пропорційно до кількості шумних нейронів (твердження~\ref{prop:broadcastcost}) & Криві навчання лінійної мережі: збурення нейронів повільніше за точний градієнт у стільки разів, скільки вихідних нейронів, збурення ваг --- ще в стільки разів, скільки входів. & \cite{Werfel2005} & теорія \\
14 & Виходи \nt{DOPA} ідуть насамперед до областей вибору дій & Повна реконструкція аксонів окремих нігростріатних \nt{DOPA}-нейронів щура: гілки однієї клітини покривають $0.45$--$5.7\,\%$ об'єму стріатума (у середньому $2.7\,\%$), поза стріатумом гілок майже немає. & \cite{Matsuda2009} & виміряно \\
15 & Кора навчається передбачати свій вхід, помилка рахується на місці & Огляд: у сенсорній корі є відповіді на розбіжність очікуваного і того входу, що надійшов. Що це загальне правило навчання кори, не доведено. & \cite{Keller2018} & гіпотеза \\
16 & Дофамін поводиться як помилка~\eqref{eq:ctd}: сплеск, перенесення на провісник, провал (рис.~\ref{fig:ctdcases}) & \nt{DOPA}-нейрони мавпи: сплеск на несподіваний сік; після навчання --- сплеск на провісник і немає відповіді на обіцяний сік; провал, коли сік не надійшов. Неперервна форма~\eqref{eq:ctd} --- теорія. & \cite{Schultz1997, Doya2000} & виміряно \\
17 & Схема, що обчислює $\dd V/\dd t$ і віднімає очікування & Миші, запис і оптогенетика у вентральній покришці: \nt{DOPA}-нейрони віднімають очікування від винагороди; сусідні \nt{GABA}-нейрони гальмують їх, коли винагорода очікується, і їхнє збудження чи гальмування змінює помилку. Як обчислюється похідна, не показано. & \cite{Eshel2015arithmetic} & гіпотеза \\
18 & \nt{DOPA}-нейрон --- ритмоводій; провали мілкіші за сплески & У зрізі \nt{DOPA}-нейрони розряджаються спонтанно і дуже регулярно, на повільному пейсмекерному потенціалі. У мавпи частота кількісно йде за додатною помилкою, а від'ємну передає лише слабко. & \cite{GraceOnn1989, Bayer2005} & виміряно \\
19 & Актор і критик (рис.~\ref{fig:actorcritic}) & Алгоритм --- з теорії навчання з підкріпленням. У людей (фМРТ) черевний стріатум поводиться як критик і під час вибору, і без нього, спинний --- лише коли треба вибирати дії. & \cite{SuttonBarto2018, ODoherty2004actor} & гіпотеза \\
20 & Знак $\delta$ у правилі перевернутий у нейронів із другою родиною рецепторів & Зрізи стріатума: у нейронів із D1-рецепторами спарювання дає посилення, якому потрібен D1; у нейронів із D2 --- послаблення, якому потрібен D2. & \cite{Shen2008} & виміряно \\
\bottomrule
\end{longtable}}
